\documentclass[11pt]{article}
\usepackage[margin=0.9in]{geometry}
\usepackage{fontspec}
\usepackage{amsmath,amssymb,amsthm}
\usepackage{xurl}
\begin{document}\sloppy
\begin{center}
{\Large GROK H0445: cube 29 of \texttt{a40ff54e580f5d0a} reaches 203 of 210 UNSAT; certified count rises to 206}\\[4pt]
2026-10-06 04:45 America/Los\_Angeles \quad branch \texttt{main}, tip \texttt{c7fdad2e} \quad automation 41774ea2\\[2pt]
\textbf{P versus NP: OPEN.}
\end{center}

\section*{Peer delta since H0427 (\texttt{c7fdad2e})}
Checked 04:46\,PT. No commits landed after H0427; \texttt{c7fdad2e} is still the tip of \texttt{main}. Nothing new under \path{research/main/} or \path{research/agents/grok/}. Root \texttt{main.tex} is not touched.

\section*{1. Literature}
None needed. The next claim is about our own finite search.

\section*{2. Claims recorded}
\textbf{Cube 29, $(0,1,2,14)$, partial (COMPUTATIONAL, solver-trusted).} The splitter log holds 203 distinct UNSAT subcube indices for cube 29, all 203 with the same step-2 tuple as the matching cube-0 subcube, no SAT line and no \texttt{TIMEOUT} line. Missing: 194, 204, 205, 206, 207, 208, 209; the first five were running at 04:46\,PT under the 600\,s cap. Slowest so far is subcube 134 at 302.2\,s, mean 43.0\,s. Cube 29 is not closed.

\textbf{Claim DRAT206 (COMPUTATIONAL).} 206 of cube 0's 210 subcubes carry drat-trim-verified refutations (new: 209, 99, 133); 4 remain; the DRAT job for subcube 52 started 04:44\,PT.

\textbf{Status of first-step cubes (COMPUTATIONAL, solver-trusted).} Closed: 0, 1, 2, 4, 9, 14, 19, 21, 22, 24, 27, so 53 of 75 kept first-step cubes are closed, 22 open, 0 SAT. Unchanged from H0427.

\section*{3. Smallest honest test}
\begin{itemize}
\item \textbf{Closure check} (\path{/workspace/scratch_h0445_1006/closure_run.py}, a copy of H0348's \path{closure.py}; output \path{closure_0446.out}). It ignores \texttt{CUBE-DONE} and needs 210 distinct UNSAT indices with matching tuples. No SAT line; the only timeout is cube 22 subcube 204, superseded by its 3000\,s rerun UNSAT (674.5\,s); closed set as above; partial \{29: 203 UNSAT, no open timeout\}.
\item \textbf{Cube scan} (\path{cscan.py}, output \path{cscan_0446.out}). Cube 29: 203 distinct, 203 of 203 tuples match, 7 missing.
\item \textbf{Recount} (\path{/workspace/scratch_h2333_1005/experiments/recount.py}, rerun 04:46\,PT). Total verified 206, remaining 4, 125 distinct CNF hashes among 125 batch lines, none not verified.
\end{itemize}

\section*{4. What this changes}
Certified rises from 203 to 206 of 210; closed stays 53 of 75. Neither moves a charged resource. The bounds stay $9\le L(\texttt{a40ff54e580f5d0a},2)\le 11$ (COMPUTATIONAL for the upper bound; the lower bound stays CONDITIONAL on the tightness lemma as recorded earlier); $L=9$ itself stays open while 22 cubes are open. The largest confirmed price stays one.

\section*{Effect on the standing candidate}
No claim is falsified and no charged resource improves. The time-state-representation surface stays CANDIDATE. The SAT-to-charged-model compiler is OPEN, and so is the product-span obstruction. These results concern one fixed function and imply no lower bound for any family.

\section*{Next smallest step}
Let cube 29 finish its last 7 subcubes and run the closure check; rerun any timeout at 3000\,s. Let the DRAT batch finish its last 4.

\bigskip
\noindent\textbf{Status sentence.} P versus NP remains OPEN: the SAT-to-charged-model compiler and the product-span obstruction are open load-bearing arrows. This pass records that cube 29 of \texttt{a40ff54e580f5d0a} at tight $L=9$ has 203 of 210 subcubes UNSAT with no timeout and is not yet closed (COMPUTATIONAL, solver-trusted), and that 206 of cube 0's 210 subcubes carry drat-trim-verified refutations (COMPUTATIONAL); 53 of 75 kept first-step cubes are closed, 22 open, 0 SAT, and $9\le L(f,2)\le 11$ is unchanged. No claim is falsified, no charged resource improves, and this is not a proof of anything about P versus NP.
\end{document}
